\documentclass[final]{beamer}
\usepackage[orientation=landscape,size=custom,width=120,height=70,scale=1.45]{beamerposter}
\usepackage[T1]{fontenc}
\usepackage[utf8]{inputenc}
\usepackage[brazilian]{babel}
\usepackage{microtype}
\usepackage{booktabs}
\usepackage{url}

\usetheme{default}
\setbeamertemplate{navigation symbols}{}
\setbeamercolor{headline}{fg=white,bg=blue!55!black}
\setbeamercolor{block title}{fg=white,bg=blue!55!black}
\setbeamercolor{block body}{fg=black,bg=blue!3}
\setbeamerfont{block title}{size=\Large}
\setbeamerfont{block body}{size=\large}

\title{Instruções de segurança para agentes de codificação com IA em software de saúde mental}
\author{Aloisio Costa}
\institute{Afiliação institucional a confirmar}
\date{Protocolo em elaboração · sem resultados empíricos}

\begin{document}
\begin{frame}[t]
\begin{block}{}
\centering
{\Huge\bfseries \inserttitle\par}
\vspace{0.45em}
{\Large \insertauthor \quad|\quad \insertinstitute\par}
\vspace{0.25em}
{\large \insertdate\par}
\end{block}
\vspace{0.7em}
\begin{columns}[t,totalwidth=\linewidth]
\begin{column}{.31\linewidth}
\begin{block}{Problema e pergunta}
Agentes de codificação alteram código e testes a partir de instruções no repositório. Este projeto examinará patches produzidos em tarefas sintéticas de engenharia associadas a um protótipo acadêmico, sem avaliar pessoas ou cuidados clínicos.

\vspace{0.5em}
\textbf{Pergunta:} como uma skill de segurança versionada altera a proporção de execuções sem violação crítica de escopo não clínico, privacidade e supervisão humana?

\vspace{0.5em}
Estudos próximos avaliam agentes conversacionais e benchmarks de respostas em saúde mental. Este protocolo observará decisões de implementação de agentes de codificação; o hiato permanece provisório até a revisão de escopo [Feng et al., 2025; Badawi et al., 2026].
\end{block}

\begin{block}{Objeto e unidade de análise}
\begin{itemize}
  \item Objeto: patches e eventos de ferramenta gerados por configurações de agentes.
  \item Unidade: uma execução agente × tarefa × condição × repetição.
  \item Sistemas: ChatGPT e Codex como configurações completas, se registráveis.
  \item Ambiente: commit limpo, tarefas inventadas, sem rede ou dados reais.
\end{itemize}
\end{block}
\end{column}

\begin{column}{.36\linewidth}
\begin{block}{Desenho planejado}
Comparação pareada por tarefa:
\begin{enumerate}
  \item Instruções base do repositório.
  \item Mesmo contexto acrescido da skill Bipolaris.
\end{enumerate}
O banco inicial prevê oito tarefas sintéticas e três repetições por célula. O total final depende do piloto e será congelado antes da fase confirmatória; esta contagem não é cálculo de poder.

\vspace{0.5em}
\textbf{Desfecho primário:} proporção de execuções com ao menos uma violação crítica, conforme rubrica pré-definida.

\textbf{Secundários:} critérios funcionais permitidos, testes, severidade, recusa excessiva e variação entre repetições.

\vspace{0.5em}
Dois revisores humanos avaliarão patches e logs. Relatar concordância, incerteza, achados nulos e ameaças à validade [Ralph et al., 2020].
\end{block}

\begin{block}{Etapas de pesquisa}
\centering
\large
Revisão de escopo \quad $\rightarrow$ \quad Piloto \quad $\rightarrow$ \quad Protocolo congelado \quad $\rightarrow$ \quad Experimento \quad $\rightarrow$ \quad Análise e relato
\end{block}
\end{column}

\begin{column}{.29\linewidth}
\begin{block}{Salvaguardas e limites}
\begin{itemize}
  \item Sem participantes, prontuários ou dados reais de saúde.
  \item Sem diagnóstico, predição, triagem, tratamento ou manejo de crise.
  \item Sem integração de modelo ao protótipo.
  \item Registrar versões, prompts, permissões, commits e testes.
  \item Qualquer coleta futura com pessoas exige determinação ética institucional antes de iniciar.
\end{itemize}
O estudo examina conformidade de engenharia em tarefas sintéticas; não mede eficácia clínica [WHO, 2024; NIST, 2024].
\end{block}

\begin{block}{Estado e próximos passos}
\textbf{Estado:} protocolo e protótipo iniciais. Revisão não concluída. Nenhum experimento ou resultado foi produzido.

\vspace{0.4em}
\textbf{Próximos gates:} concluir revisão; validar tarefas e rubrica; pilotar; documentar decisão ética; congelar análise; executar estudo; atualizar artigo e banner com achados observados.
\end{block}

\begin{block}{Projeto e referências}
Repositório: \url{https://github.com/aloisiocosta-prof/bipolaris}

\vspace{0.4em}
\scriptsize
Feng et al. (2025). \textit{Journal of Medical Internet Research}, 27, e69639. DOI: 10.2196/69639.\par
Badawi et al. (2026). \textit{EACL}, p. 3873--3896. DOI: 10.18653/v1/2026.eacl-long.180.\par
Ralph et al. (2020). \textit{Empirical Standards for Software Engineering Research}.\par
WHO (2024). \textit{Ethics and governance of AI for health}. NIST (2024). \textit{Generative AI Profile}, AI 600-1.
\end{block}
\end{column}
\end{columns}
\end{frame}
\end{document}
